\documentclass[carta,firma,sinnotas]{lafrox}
\usepackage{fvextra}
\usepackage{pdflscape}

\datoslafrox{
  documento = {Formulario T-16: Plan de riesgos},
  subtitulo = {Registro de las amenazas},
  alcance = {Gestión de riesgos de la propuesta LafroX},
  formulario = {T-16: Plan de riesgos},
  fecha = {27 de Octubre de 2026},
  version = {0.1},
  nomenclatura = {LAFROX-Formulario-T-16},
}
% Soporte local de firma para el cuerpo y los anexos del Subdocumento 8.
% No altera la clase ni las plantillas compartidas.
\newcommand{\LFXportadasd}{%
  \AddToHookNext{shipout/foreground}{%
    \begin{tikzpicture}[remember picture,overlay]
      \node[anchor=north west,inner sep=0pt]
        at \LFXen{\LFXmi}{0.835}
        {\includegraphics[height=13mm]{logo/FIRMA_Alex_Aravena_transparente.png}};
    \end{tikzpicture}%
  }%
  \portadalafrox
}

% Firma abreviada sobre filete, en el mismo espacio del pie existente.
\renewcommand{\LFXpiederecho}{%
  \begingroup\LFXpiebase
  \begin{tabular}[b]{@{}r@{}}
    \raisebox{0pt}[0pt][0pt]{%
      \begin{tabular}[b]{@{}c@{}}
        \includegraphics[height=7mm]{logo/FIRMA_Alex_Aravena_transparente.png}\\[-1pt]
        \rule{3.4cm}{0.4pt}
      \end{tabular}%
    }\\[-1pt]
    {\intermedium\fontsize{9pt}{11pt}\selectfont\color{lafroxTinta}%
      Media firma: representante\quad\LFXfoliotexto}
  \end{tabular}%
  \endgroup
}


% Formato de portada común, basado en el Subdocumento 1.
\configurarportadalafrox{Formulario T-16 \textemdash{} Plan de riesgos}{Subdocumento 08}{Formulario T-16}{LAFROX-Formulario-T-16.pdf}{logo/FIRMA_Alex_Aravena_transparente.png}
\begin{document}
\portadalafrox
\preliminareslafrox
\setcounter{chapter}{7}

% Fuente Markdown: LAFROX-Formulario-T-16.md
% Contenido actualizado; estructura y formato lafrox.cls conservados.

El registro sigue el formato del Formulario T-16, incluido en las Bases Administrativas (Bases Administrativas, Formulario T-16, p. 63), y la norma ISO 31000 (International Organization for Standardization [ISO], 2018), conforme al RT-19.04 (Bases Técnicas Transversales, cap. 19, RT-19.04, p. 33). Sus riesgos se anclan en el Caso 02 (Bases Técnicas del caso, caps. 10 y 19, pp. 19--20, 35--36), cubren los temas del índice obligatorio indicado en el Formulario T-16 (Bases Administrativas, Formulario T-16, p. 63) y se desarrollan en el Subdocumento 8. El registro contiene los 35 riesgos declarados al CLIENTE; conserva los ID originales con discontinuidades. Los 16 asuntos internos retirados, incluidos los 13 indicados por la revisión, se siguen como CI en el Anexo 8.E y no integran esta tabla. Prob. presenta el nivel cualitativo de probabilidad y su factor; Impacto presenta el nivel global cualitativo y su factor. El nivel de probabilidad se deriva de $p_0$ del Anexo 8.B mediante las bandas de la Tabla 8.2; su factor es una ponderación de categoría. El impacto es el mayor nivel por esfuerzo condicional, plazo, alcance, calidad sanitaria o consecuencias contractuales (Tabla 8.3 y respaldo por ID en el Anexo 8.C). Exposición es $q_P\times q_I$ y la zona resulta de los umbrales de la Tabla 8.4. La columna Mitigación presenta la estrategia y la respuesta propuesta; el disparador observable se detalla en la Tabla 8.A.1. Para los 28 riesgos con coordinación del CLIENTE se explicitan contraparte, solicitud, plazo y contingencia, con evidencia requerida en la Tabla 8.A.2. CT significa Contraparte Técnica; D es la fecha de la actividad indicada y D-10/D-5 son días hábiles anteriores a ella. El protocolo de solicitud y escalamiento se define en 8.1.2 del Subdocumento 8; los plazos de coordinación son propuestos y no alteran la recepción formal de entregables. El responsable corresponde al dueño propuesto en el análisis. Estos juicios cualitativos no son frecuencias observadas ni costos monetarios. El Anexo 8.D vincula las brechas con CI-048, R-049, CI-050 y R-042 y verifica las demás omisiones. CI-048 y CI-050 son verificaciones internas, no riesgos transferidos al CLIENTE. El Anexo 8.B explicita las hipótesis de probabilidad, episodios y HH de la ampliación. El detalle de valores esperados está en el Anexo 8.B; la evaluación de respuestas se desarrolla en el cuerpo del Subdocumento 8. Cada responsable dirige su equipo. DAT e IMP conservan su familia de esfuerzo, pero desde el mes 22 la responsabilidad pasa al Líder de Operación / SRE (SD8, Tabla 8.1). Los ID retirados no se reutilizan ni se renumeran. No se establece aquí una reserva cuantitativa.

\begin{landscape}
\setlength{\textwidth}{\linewidth}
\begingroup
\renewcommand{\thetable}{T-16.1}
% Compactación local: conservar legibilidad y columnas del formulario, reduciendo
% el espacio vertical para que cada riesgo no ocupe una página completa.
\renewcommand{\LFXcuerpotabla}{\fontsize{7.5pt}{8.5pt}\selectfont\addfontfeature{}}
\renewcommand{\arraystretch}{1.05}
\setlength{\tabcolsep}{2pt}

\begin{tablalafrox}{Registro cualitativo de riesgos; Fuente: contenido de licitación y subdocumentos anteriores}{tab:t16-1}{G{0.05} Y F{0.17} G{0.10} G{0.06} G{0.05} Y F{0.12}}{\cab{ID} & \cab{Riesgo} & \cab{Categoría} & \cab{Prob.} & \cab{Impacto} & \cab{Expos.} & \cab{Mitigación} & \cab{Responsable}}
R-001 & Dado que la habilitación EDI de cada cadena depende de su participación y certificación externa antes de las cuatro semanas finales de evidencia de E2, entonces puede que una cadena incluida en el alcance no complete su participación o certificación EDI antes de la fecha requerida, teniendo como consecuencia que no se acredite completo el intercambio electrónico exigido en E2 y deban reprogramarse pruebas y coordinación. & Integración y legado & 3 / 0,50 & 4 / 0,40 & 0,20 (Alta) & Solicitar al Gerente Comercial, con TI y cada cadena, agenda EDI en mes 13, confirmada en mes 14, y certificación a D-5 del inicio de las cuatro semanas finales de evidencia de E2. Ante falta de respuesta, JP escala por CT; continuar pruebas independientes sin acreditar la cadena pendiente ni sustituir EDI por carga manual. & Arquitecto de Solución \\
R-002 & Dado que la migración sanitaria utiliza antecedentes del ERP, WMS y planillas y requiere vínculos históricos verificables con evidencia de origen, entonces puede que los datos recibidos contengan faltantes de lote o vínculos históricos que no puedan resolverse con esa evidencia, teniendo como consecuencia que la carga sanitaria afectada quede pendiente de conciliación y decisión documentada, y aumente el trabajo de saneamiento. & Solución y tecnología & 4 / 0,70 & 5 / 0,80 & 0,56 (Crítica) & Solicitar a la Jefa de Calidad, con TI y Bodega, evidencia de origen y decisión sobre lotes no reconstruibles; política V-07 al cierre del mes 2 y excepciones a D-5 de validar cada carga. Si no responde, JP escala por CT, se aísla el conjunto y queda pendiente su aceptación; no se inventa historia sanitaria. & Líder de Datos \\
R-003 & Dado que la migración y la reposición dependen de la extensión y calidad efectiva del historial de ventas disponible en el ERP, entonces puede que el historial extraído tenga períodos incompletos o datos que no puedan utilizarse directamente, teniendo como consecuencia que deban ampliarse la preparación y conciliación de la carga y el tratamiento de excepciones de reposición. & Integración y legado & 3 / 0,50 & 3 / 0,20 & 0,10 (Media) & Solicitar al Jefe de Abastecimiento y a TI el historial de tres años y validación de V-09 al cierre del mes 3; corregir extractos a D-5 de validar la carga. Ante silencio, escalar por CT, separar períodos incompletos y continuar el perfilado disponible sin reducir la historia exigida. & Líder de Datos \\
R-004 & Dado que la aceptación del inventario migrado requiere conciliar los saldos legados con el conteo físico y aprobar las explicaciones y ajustes antes del corte, entonces puede que se detecten diferencias sin explicación y ajuste aprobados a tiempo, teniendo como consecuencia que deban repetirse el saneamiento y la validación, y la carga afectada no pueda ser aceptada hasta resolverlas. & Operación del cliente & 4 / 0,70 & 4 / 0,40 & 0,28 (Alta) & Solicitar a Bodega, con Operaciones y Finanzas, protocolo de conciliación en mes 2, validado en mes 3. Confirmar conteo a D-5 del corte y cerrar acta y ajustes antes de D. Si no hay respuesta, escalar por CT, repetir conciliación y bloquear el corte afectado hasta contar con evidencia aprobada. & Líder de Datos \\
R-005 & Dado que la carga inicial de M8 requiere saldos de envases por cliente y transportista sustentados en antecedentes conciliables del cliente, entonces puede que los antecedentes entregados no permitan determinar un saldo conciliado por cliente y transportista, teniendo como consecuencia que deban repetirse la preparación y validación de la carga y no pueda demostrarse M8 con saldos iniciales aceptados. & Operación del cliente & 3 / 0,50 & 4 / 0,40 & 0,20 (Alta) & Solicitar al Gerente de Operaciones y a Finanzas respaldos y conciliación V-08 de saldos por cliente y transportista; política en mes 2 y respuesta a D-5 de validar la carga M8 de cada ola. Ante silencio, escalar por CT y mantener M8 pendiente para esos saldos; no presumir saldo cero. & Líder de Datos \\
R-006 & Dado que las comparaciones históricas de aceptación dependen de horas verificables de entrega y montos esperados de rendición que la captura futura no reconstruye, entonces puede que falten horas verificables de entrega o montos esperados de rendición en los antecedentes históricos, teniendo como consecuencia que no puedan respaldarse con esos registros las comparaciones de OTIF o cuadratura y deba repetirse la validación del período de referencia. & Operación del cliente & 4 / 0,70 & 4 / 0,40 & 0,28 (Alta) & Solicitar a Comercial y Operaciones las horas históricas y a Finanzas los montos de rendición, con validación de V-10 al cierre del mes 1. Si no responden, escalar por CT y mantener pendiente la referencia de aceptación; la captura futura no reconstruye la historia anterior. & Líder de Implantación y Gestión del Cambio \\
R-007 & Dado que el volumen histórico utilizado para estimar la carga de cartera no acredita que incluya todos los documentos de deuda todavía vigente, entonces puede que se identifiquen documentos de deuda vigente fuera del volumen histórico considerado, teniendo como consecuencia que deban extraerse, transformarse y conciliarse datos adicionales, aumentando el trabajo y el volumen de transferencia de la migración de cartera. & Integración y legado & 3 / 0,50 & 3 / 0,20 & 0,10 (Media) & Solicitar a Finanzas y TI el universo de deuda vigente, con respuesta inicial al cierre del mes 3 y conciliación a D-5 de validar la migración. Ante silencio, escalar por CT y mantener la cartera incompleta pendiente; preparar extracción complementaria sin omitir deuda obligatoria. & Líder de Datos \\
R-008 & Dado que los trámites tributarios se envían por el sistema de gestión actual del cliente y requieren respuesta del Servicio de Impuestos Internos (SII) para conocer su resultado, entonces puede que el SII no responda mientras se tramita un documento, teniendo como consecuencia que el trámite quede sin confirmación y se retrase su cierre mientras se gestiona y comprueba la respuesta, sin que el silencio equivalga a aprobación. & Integración y legado & 1 / 0,10 & 3 / 0,20 & 0,02 (Baja) & Aceptar pasivamente. Mantener la observación de pendientes y los controles ya previstos para consultar y conciliar el resultado tributario. No proponer inversión preventiva adicional con esta valoración. La ausencia de respuesta no autoriza aprobar el trámite ni cerrar el proceso. & Líder de Operación / SRE \\
R-011 & Dado que la certificación de la parametrización térmica requiere una matriz de umbrales y duraciones por producto aprobada por Calidad del cliente, entonces puede que Calidad no apruebe esa matriz antes del cierre de las pruebas y la marcha blanca de los productos afectados, teniendo como consecuencia que no pueda completarse la evidencia de parametrización y aceptación del alcance térmico y deban reprogramarse validaciones. & Operación del cliente & 3 / 0,50 & 4 / 0,40 & 0,20 (Alta) & Solicitar a la Jefa de Calidad la matriz térmica versionada V-11 durante mes 1 y su aprobación al cierre del mes 2, antes de certificar. Si no responde, escalar por CT y mantener el tratamiento preventivo de lecturas fuera de ficha; no cerrar marcha blanca ni aprobar parámetros por el CLIENTE. & Líder de Calidad \\
R-012 & Dado que el cliente compra el equipamiento de terreno conforme a la especificación de LafroX y debe disponer de él antes de cada ola, entonces puede que el cliente o su proveedor no entregue el equipamiento antes de la ola programada, teniendo como consecuencia que no puedan ejecutarse con esos equipos la instalación y las pruebas previstas y deban reprogramarse despliegue y acompañamiento. & Operación del cliente & 2 / 0,30 & 4 / 0,40 & 0,12 (Media) & Solicitar al Jefe de TI, con Operaciones, calendario de equipos en mes 1, entrega confirmada a D-5 y recepción técnica antes de la instalación o prueba D. Ante silencio o faltante, escalar por CT, continuar trabajo independiente y bloquear las actividades sin equipo. & Líder de Implantación y Gestión del Cambio \\
R-013 & Dado que la incorporación de cada empresa transportista al reparto requiere que su acuerdo operativo esté firmado antes de la ola correspondiente, entonces puede que una empresa transportista no firme el acuerdo a tiempo, teniendo como consecuencia que sus conductores no puedan incorporarse ni completar las pruebas de esa ola y quede pendiente demostrar la cobertura del transportista. & Personas y adopción & 3 / 0,50 & 4 / 0,40 & 0,20 (Alta) & Solicitar al Gerente de Operaciones y al representante de cada transportista el acuerdo y la nómina desde mes 2, firmados a D-5 de su incorporación a la ola. Ante silencio, escalar por CT y no habilitar a esa empresa; conservar su cobertura pendiente sin retirarla del alcance. & Líder de Implantación y Gestión del Cambio \\
R-014 & Dado que la instalación de termógrafos en camiones refrigerados de terceros requiere acceso al vehículo y autorización de su propietario, entonces puede que el propietario niegue el acceso o la autorización para instalar el equipo, teniendo como consecuencia que no pueda instalarse ni integrarse el termógrafo antes de incorporar la ruta y quede pendiente la evidencia de temperatura. & Operación del cliente & 2 / 0,30 & 4 / 0,40 & 0,12 (Media) & Solicitar a Operaciones y al propietario permiso de instalación por vehículo con D-10 de la visita y obtenerlo a D-5. Ante silencio, escalar por CT, suspender la intervención y proponer otra visita o vehículo autorizado; mantener pendiente la evidencia térmica. & Líder de Implantación y Gestión del Cambio \\
R-015 & Dado que la implantación exige que los usuarios incorporen sus transacciones al registro oficial de la solución, entonces puede que parte de los usuarios continúe registrando exclusivamente en el procedimiento anterior durante una ola, teniendo como consecuencia que no pueda acreditarse la cobertura de adopción requerida para cerrar esa ola y deba extenderse la intervención de implantación. & Personas y adopción & 4 / 0,70 & 4 / 0,40 & 0,28 (Alta) & Solicitar a Comercial y Operaciones, con Bodega por turno, responsables y agenda de acompañamiento a D-5 del inicio de cada ola, convocando a D-10. Si no responden o persiste el registro externo, escalar por CT, focalizar acompañamiento y mantener pendiente la evidencia de adopción. & Líder de Implantación y Gestión del Cambio \\
R-016 & Dado que el planificador del CLIENTE se jubila dentro del horizonte contractual y su ausencia reduce el cumplimiento de entrega entre 6 y 9 puntos porcentuales según el caso, entonces puede que llegue su retiro o indisponibilidad antes de transferir y probar las reglas y excepciones de rutas, teniendo como consecuencia pérdida de conocimiento operativo, reconstrucción extraordinaria y riesgo de no acreditar el resultado R18-16. & Personas y adopción & 2 / 0,30 & 4 / 0,40 & 0,12 (Media) & Solicitar al Gerente de Operaciones y al planificador confirmación de retiro y disponibilidad V-14 al cierre del mes 1; transferir y validar reglas con el ayudante en meses 1--3. Probar las rutas y excepciones sin el planificador en la marcha blanca de meses 13--15 (S7, 7.2.1). Si falta respuesta o se retira antes, escalar por CT y reconstruir con ayudante e historia verificable, sin dar por validadas reglas pendientes. & Líder de Implantación y Gestión del Cambio \\
R-017 & Dado que la capacitación y certificación deben realizarse sin detener pedidos, ventas ni reparto, y el cliente debe disponer de los usuarios en los horarios acordados, entonces puede que las necesidades operativas impidan que algunos usuarios asistan a formación y evaluación antes del cierre de su ola, teniendo como consecuencia que deban reprogramarse sesiones y acompañamiento y no pueda acreditarse que esos usuarios están preparados hasta que completen la evaluación. & Personas y adopción & 2 / 0,30 & 4 / 0,40 & 0,12 (Media) & Solicitar a Operaciones, Bodega, Comercial y TI, según perfil, nómina y liberación horaria a D-5 de cada sesión convocada a D-10. Ante silencio o ausencia, escalar por CT y reprogramar evaluación; usar sólo suplentes certificados y no acreditar preparación ni cierre de ola sin evidencia. & Líder de Implantación y Gestión del Cambio \\
R-018 & Dado que los diseños, pruebas y aceptaciones incluyen actividades que dependen de validaciones o decisiones de áreas del cliente, entonces puede que el área requerida no entregue una validación o decisión en la fecha acordada, teniendo como consecuencia que se detengan o repitan tareas dependientes, se consuma holgura y deban reprogramarse recursos. & Operación del cliente & 3 / 0,50 & 4 / 0,40 & 0,20 (Alta) & Solicitar a CT área decisora y suplente, con agenda al cierre del mes 1; cursar insumos a D-10 y pedir respuesta a D-5. Las aceptaciones formales conservan diez días hábiles. Para promover la región de contingencia, avisar al suplente a los 15 min sin respuesta del titular y obtener decisión dentro de 30 min (S4, 4.3.2.5); no esperar el plazo ordinario. Mantener bloqueada la acción sin autorización; el silencio no permite cambios. & Jefe de Proyecto \\
R-019 & Dado que Concepción parte con fibra de un solo proveedor sin respaldo y los demás sitios también dependen de enlaces externos, entonces puede que se interrumpa el camino disponible antes de verificar los respaldos propuestos o falle el último camino utilizable durante operación, teniendo como consecuencia aislamiento de servicios centrales, acumulación de sincronización y gestión de restitución y conciliación. & Operación del cliente & 2 / 0,30 & 4 / 0,40 & 0,12 (Media) & Mitigar. Verificar respaldos antes de aceptar la ola de Concepción; no considerar disponible un enlace sólo por estar ofertado. Mantener y ensayar los enlaces ya ofertados: en Talca y Concepción, fibra y LTE con Starlink en espera caliente; en cross-docking, Starlink con LTE Entel/Claro y conmutación menor a 30 s. En Los Ángeles de madrugada, probar pérdida del camino Starlink restante y continuidad local de su ventana de 3 h. Coordinar restitución y conciliación sin duplicar el mismo corte como falla de SII, POS o mapas. Referencia: S4, 4.2.5.3. & Líder de Operación / SRE \\
R-020 & Dado que algunos equipos de infraestructura dependen del suministro eléctrico y tienen respaldo de autonomía limitada, entonces puede que un corte de electricidad supere la autonomía del respaldo aplicable en un sitio atendido por LafroX, teniendo como consecuencia que los equipos afectados dejen de prestar servicio hasta restablecer la energía o recuperar la infraestructura. & Operación del cliente & 1 / 0,10 & 4 / 0,40 & 0,04 (Baja) & Aceptar pasivamente. Observar suministro y autonomía del respaldo existente. Si se supera, atender solo los componentes mantenidos por LafroX y verificar su restitución. Conservar la zona baja sin añadir capacidad eléctrica ni atribuir pérdidas operacionales del cliente al impacto del proyecto. & Líder de Operación / SRE \\
R-021 & Dado que la solución propone operar y recuperarse en regiones de nube de Brasil y Estados Unidos, y el tratamiento y transferencia internacional requieren aprobación del cliente, entonces puede que el cliente no otorgue la aprobación necesaria sobre residencia y transferencia de datos al habilitar esos entornos, teniendo como consecuencia que no pueda completarse la habilitación o la evidencia de recuperación afectada hasta resolver formalmente la condición. & Regulatorio y externo & 2 / 0,30 & 4 / 0,40 & 0,12 (Media) & Solicitar por CT, con TI y la autoridad del CLIENTE encauzada por Gerencia General, aprobación escrita de residencia y transferencia en mes 1, resuelta al cierre del mes 2 o antes de habilitar datos, lo que ocurra primero. Ante silencio, escalar y preparar alternativa de arquitectura; mantener deshabilitada la transferencia no aprobada. & Arquitecto de Solución \\
R-022 & Dado que una intervención para remediar un incidente o vulnerabilidad durante un período congelado requiere autorización del cliente, entonces puede que el cliente no autorice el cambio de emergencia solicitado, teniendo como consecuencia que deba sostenerse la contención y gestionarse el conflicto entre la fecha habilitada y el plazo de remediación, sin que ello implique una exención contractual. & Operación del cliente & 2 / 0,30 & 4 / 0,40 & 0,12 (Media) & Solicitar a CT, Operaciones y TI flujo y suplentes al cierre del mes 1; ante incidente, pedir autorización escrita antes del cambio. Sin respuesta, escalar en una hora o antes si lo exige el plazo de atención, manteniendo contención autorizada y estado seguro. El silencio no permite intervenir durante congelamiento. & Encargado de Seguridad de la Información \\
R-028 & Dado que los recorridos de prueba se realizan en camiones cuyo funcionamiento mecánico es necesario para completarlos, entonces puede que un camión sufra una falla mecánica durante el recorrido, teniendo como consecuencia que se interrumpa el recorrido y no se completen los casos de prueba programados, por lo que deba repetirse la validación. & Operación del cliente & 2 / 0,30 & 2 / 0,10 & 0,03 (Baja) & Aceptar pasivamente. Observar la disponibilidad del recorrido y, si una falla interrumpe las pruebas, coordinar su repetición con el cliente. La zona baja no justifica por sí sola una acción preventiva nueva. Limitar el efecto de LafroX a tiempo y recursos de validación, sin cubrir mantenimiento del camión. & Líder de Implantación y Gestión del Cambio \\
R-029 & Dado que los casos de prueba y la evidencia de aceptación requieren completar el recorrido del camión, entonces puede que el camión se vea envuelto en un accidente de tránsito durante el recorrido, teniendo como consecuencia que se interrumpan los casos de prueba y deba coordinarse la repetición de validaciones, consumiendo tiempo y recursos de implantación. & Operación del cliente & 1 / 0,10 & 2 / 0,10 & 0,01 (Baja) & Aceptar pasivamente. Ante aviso de siniestro, coordinar con el cliente otro recorrido autorizado para completar los casos pendientes. Mantener vigilancia ordinaria sin nueva respuesta preventiva. No imputar a LafroX lesiones, indemnizaciones ni reposición de vehículo o carga. & Líder de Implantación y Gestión del Cambio \\
R-030 & Dado que los recorridos de prueba requieren que el camión llegue al destino previsto para completar los casos y reunir evidencia, entonces puede que el camión sea asaltado, robado o incendiado durante el recorrido, teniendo como consecuencia que se interrumpa la validación y deban reprogramarse y repetirse las pruebas afectadas, con posible consumo adicional de tiempo y recursos. & Operación del cliente & 1 / 0,10 & 2 / 0,10 & 0,01 (Baja) & Aceptar pasivamente. Ante interrupción por asalto, esperar la habilitación segura del cliente y coordinar repetición de las validaciones pendientes. La aceptación cubre la exposición del proyecto en zona baja. No asumir custodia de bienes, pérdidas delictuales o responsabilidad por seguridad del transporte. & Líder de Implantación y Gestión del Cambio \\
R-031 & Dado que los conductores externos requieren comunicación, autorización y capacitación para utilizar la solución, entonces puede que un transportista presente sin aviso a un reemplazo que aún no esté autorizado o capacitado, teniendo como consecuencia que deban atenderse su habilitación y formación y verificarse las entregas afectadas antes de cerrar la evidencia de la ola. & Personas y adopción & 3 / 0,50 & 3 / 0,20 & 0,10 (Media) & Solicitar a Operaciones y al transportista nómina y suplentes a D-5 de incorporar la empresa y confirmar conductor, vehículo y viaje antes de cada salida. Si falta respuesta o llega un sustituto, escalar, verificar en andén y solicitar conductor certificado; no habilitar uso no autorizado ni acreditar cobertura incompleta. & Líder de Implantación y Gestión del Cambio \\
R-032 & Dado que la solución tiene un alcance delimitado por requisitos, exclusiones y decisiones aprobadas, contra el cual deben evaluarse las solicitudes posteriores, entonces puede que la contraparte solicite una función, integración o modificación que exceda el alcance y requiera trabajo adicional, teniendo como consecuencia que deba evaluarse y tramitarse el cambio y, si se aprueba, ajustarse esfuerzo, costo o programación según el impacto acordado. & Contrato y licitación & 3 / 0,50 & 3 / 0,20 & 0,10 (Media) & Solicitar por CT a la gerencia peticionaria solicitud escrita y decisión del cambio antes de su primera actividad dependiente; registrar y remitir el mismo día hábil. Ante silencio, mantener línea base y continuar trabajo independiente; si amenaza un hito, escalar al Comité Ejecutivo sin iniciar desarrollo no aprobado. & Jefe de Proyecto \\
R-033 & Dado que la solución trata y conserva documentos, ubicación y evidencia durante una vigencia contractual de varios años sujeta a exigencias normativas, entonces puede que una modificación normativa o nueva interpretación aplicable exija ajustar controles, conservación o eliminación de datos, teniendo como consecuencia que deban adaptarse y validarse componentes y procedimientos afectados, con trabajo adicional y posibles ajustes de programación sujetos al contrato. & Regulatorio y externo & 2 / 0,30 & 4 / 0,40 & 0,12 (Media) & Escalar. La interpretación normativa y la decisión sobre obligaciones y cambios contractuales requieren autoridad del cliente y asesoría competente. Elevar el cambio aplicable para definir alcance y financiación, y luego planificar adecuación técnica. No inventar reformas, multas ni cobertura contractual. & Arquitecto de Solución \\
R-039 & Dado que el tramo variable de la Operación de la innovación 4 se liga al costo de servir, entonces puede que ese costo varíe por causas externas, como el combustible, un cambio de malla o el peak de septiembre, o que se intente bajarlo degradando el servicio o excluyendo clientes, teniendo como consecuencia disputas en la liquidación del tramo variable. & Contrato y licitación & 3 / 0,50 & 3 / 0,20 & 0,10 (Media) & Solicitar a Finanzas, con Comercial y CT, reglas de 5.4.3 aprobadas al cierre del mes 20 y tres liquidaciones en sombra reproducidas en meses 21--23 antes del pago variable del mes 24. Ante silencio, no activar el tramo; una disputa se escala y suspende sólo el devengo correspondiente según contrato, manteniendo el servicio. & Jefe de Proyecto \\
R-041 & Dado que existe objeción sindical a cámaras en cabina y al control de jornada por GPS, aunque la propuesta excluye cámaras en cabina y separa la telemetría operacional, entonces puede que no se acuerden o comprendan las finalidades y controles del uso propuesto antes de la ola de reparto, teniendo como consecuencia resistencia a la incorporación y pruebas o aceptaciones pendientes. & Personas y adopción & 3 / 0,50 & 4 / 0,40 & 0,20 (Alta) & Mitigar y escalar. JP coordina con Operaciones, TI y sindicato el acuerdo 5.4.2 antes de la ola: convocar a D-10 y resolver a D-5. Explicar finalidades, acceso y controles; mantener excluidas las cámaras en cabina y no usar GPS para jornada sin acuerdo previo. Ante desacuerdo, escalar por CT, mantener sólo el uso operacional autorizado y registrar la cobertura de la ola pendiente. & Jefe de Proyecto \\
R-042 & Dado que preparación nocturna tiene rotación anual del 38\% y requiere capacitación continua, entonces puede que las incorporaciones superen la capacidad de tutoría y evaluación o lleguen al flujo sin competencias verificadas, teniendo como consecuencia errores de preparación, refuerzo extraordinario y certificación de la ola pendiente. & Personas y adopción & 4 / 0,70 & 4 / 0,40 & 0,28 (Alta) & Mitigar. Mantener 7.1.2 en meses 13--56, turno 22:00--06:00, tutor y evaluación antes del trabajo autónomo, con aprendizaje máximo de dos horas. Bodega y Operaciones entregan altas por turno antes de asignarlas. IMP conduce hasta mes 21 y SRE desde 22, familia IMP. Ante sobrecarga, usar personal certificado, reforzar la inducción y escalar por CT; no duplicar las HH ordinarias de 7.1.2. & Líder de Implantación y Gestión del Cambio \\
R-043 & Dado que el CLIENTE no tiene documentación técnica de las interfaces del ERP ni del WMS de 2013, que se reemplaza progresivamente, entonces puede que el levantamiento no permita verificar a tiempo contratos de intercambio y extracción del legado, teniendo como consecuencia reproceso de adaptadores o cargas y bloqueo de pruebas de integración o migración. & Integración y legado & 3 / 0,50 & 4 / 0,40 & 0,20 (Alta) & Mitigar y escalar. ARQ solicita a TI acceso, diccionario y condiciones de extracción en mes 1 (1.2.3, 3.3.2 y V-16); DAT verifica la migración del WMS. Probar temprano sin modificar el ERP. Si falta una interfaz, evaluar intercambio de archivos controlado, sujeto a demostrar los requisitos; elevar por CT el bloqueo y sus impactos. La alternativa no acredita por sí sola latencia, integridad ni emisión tributaria. & Arquitecto de Solución \\
R-044 & Dado que el peak de septiembre casi duplica entregas durante tres semanas y coincide con el cierre de la marcha blanca E2 en el calendario supuesto de S7, entonces puede que un incidente, una validación pendiente o la falta de disponibilidad de usuarios no pueda resolverse dentro de las ventanas autorizadas, teniendo como consecuencia incumplimiento de condiciones de cierre y necesidad de recuperación de la programación. & Contrato y licitación & 3 / 0,50 & 4 / 0,40 & 0,20 (Alta) & Mitigar y escalar. JP acuerda con Operaciones y CT disponibilidad y continuidad antes de septiembre; cerrar habilitaciones en agosto según S7. En meses 19--20 mantener evidencia diaria de las seis condiciones de marcha blanca. Un defecto durante 1--25 de septiembre no autoriza despliegue: activar continuidad aprobada y R-022 si requiere autorización. Si persiste una condición incumplida, no declarar cierre; elevar plan de recuperación sin mover hitos por cuenta propia. & Jefe de Proyecto \\
R-045 & Dado que la propuesta opera durante 56 meses con componentes como Laravel 13, PHP 8.5 y Proxmox, entonces puede que cambien su soporte o compatibilidad de forma que las actualizaciones ordinarias no basten, teniendo como consecuencia una adecuación extraordinaria y revalidación para mantener seguridad y continuidad. & Solución y tecnología & 3 / 0,50 & 4 / 0,40 & 0,20 (Alta) & Mitigar. ARQ mantiene inventario de versiones y revisa soporte y compatibilidad trimestralmente desde mes 1, con DES y SRE. Preparar y probar adecuación antes del límite verificado de soporte, dentro de ventanas autorizadas. Si el cambio excede la actualización ordinaria 8.2.5, registrar el esfuerzo extraordinario y elevar impactos por CT; usar reversión ensayada y soporte permitido, sin sostener una versión insegura por conveniencia. & Arquitecto de Solución \\
R-046 & Dado que la arquitectura utiliza servicios gestionados de AWS y capacidades específicas de su plataforma, entonces puede que un cambio de condiciones o una necesidad de salida encuentre restricciones de exportación, credenciales o equivalencia funcional no resueltas, teniendo como consecuencia trabajo extraordinario y demora para sustituir el servicio o ejecutar reversibilidad. & Solución y tecnología & 2 / 0,30 & 4 / 0,40 & 0,12 (Media) & Mitigar. ARQ verifica con TI propiedad de cuentas, acceso a datos, claves, exportaciones e infraestructura como código en la recepción inicial y en cada actualización anual del plan 8.1.4. Probar una restauración utilizable, no sólo descargar archivos. Ante restricción, documentar brecha, mantener continuidad autorizada y escalar por CT alternativas y decisión contractual; no declarar reversibilidad completa ni cambiar de proveedor sin prueba y aprobación. & Arquitecto de Solución \\
R-049 & Dado que el ERP del CLIENTE es el único emisor tributario y reside en Talca, entonces puede que su indisponibilidad, incluida la pérdida de la sala, impida emitir la nueva guía requerida por un cambio de carga o una anulación, teniendo como consecuencia bloqueo de la salida hasta restituir el ERP. & Integración y legado & 2 / 0,30 & 4 / 0,40 & 0,12 (Media) & Mitigar y escalar. ARQ/SRE acuerdan con TI y Finanzas restitución del ERP y ensayan AL-DTE-01/AL-DR-01 antes de H5/H10, incluida pérdida de sala. Mantener sólo salidas con guía preemitida aún válida y bloquear las que requieren nueva guía. Recuperar WMS en nube no restaura el ERP ni exceptúa el RTO requerido; registrar y escalar por CT cualquier incumplimiento, sin segundo emisor ni aprobación tributaria por silencio. & Arquitecto de Solución \\
R-051 & Dado que los conductores de terceros no son trabajadores de Puelche y requieren incorporación acordada con su transportista, entonces puede que un conductor autorizado y convocado rechace utilizar el terminal o el flujo digital aun existiendo acuerdo de la empresa, teniendo como consecuencia entregas sin evidencia digital completa y cobertura de adopción de la ola pendiente. & Personas y adopción & 3 / 0,50 & 4 / 0,40 & 0,20 (Alta) & Mitigar y escalar. IMP coordina con Operaciones, transportista y conductor demostración y acompañamiento antes de su salida, dentro de 7.1.3. Confirmar disponibilidad a D-5 de la ola y aceptación práctica en andén. Ante rechazo, solicitar a la empresa conductor certificado y dispuesto, mantener acompañamiento y elevar por CT la cobertura pendiente; no atribuir a LafroX facultades laborales ni declarar completa la evidencia faltante. & Líder de Implantación y Gestión del Cambio \\
\end{tablalafrox}

\endgroup
\end{landscape}

\section*{Referencias}
\addcontentsline{toc}{section}{Referencias}

Las fuentes citadas en este documento se listan en formato APA 7.ª edición. Las Bases se citan en el texto con su documento y el artículo, capítulo, sección o código del requisito.

\begingroup
\small\setlength{\parindent}{-0.5in}\setlength{\leftskip}{0.5in}\setlength{\parskip}{3pt}
Distribuidora Puelche S.A. (2026a). \emph{Bases Administrativas de Licitación N.º TFEP-01/2026: Contratación de Solución Integral de Software y Servicios de Operación}.

Distribuidora Puelche S.A. (2026b). \emph{Bases Técnicas Transversales de Licitación N.º TFEP-01/2026}.

Distribuidora Puelche S.A. (2026c). \emph{Caso 02: Logística. Especificaciones del problema y operación de Distribuidora Puelche S.A.}

International Organization for Standardization. (2018). \emph{ISO 31000:2018 Risk management \textemdash{} Guidelines}. https://www.iso.org/standard/65694.html


\endgroup




\end{document}
